\chapter*{Заключение}                       % Заголовок
\addcontentsline{toc}{chapter}{Заключение}  % Добавляем его в оглавление

В ходе проведенного исследования были рассмотрены теоретико-методологические основы анализа промышленного цикла, систематизированы основные марксистские, буржуазные и эклектические концепции его происхождения и механизма, уточнено содержание категории промышленного цикла и раскрыта ее связь с воспроизводством общественного капитала, движением нормы прибыли, накоплением, обновлением основного капитала, кредитом и кризисами перепроизводства. Кроме того, были исследованы фазовая структура и историческая эволюция промышленного цикла, особенности его проявления в современной капиталистической экономике, роль государственного регулирования и международного движения капитала, а также проведен эмпирический анализ циклической динамики на основе официальных статистических данных по экономике США.

По результатам исследования сделан ряд выводов:
\begin{enumerate}
    \item Промышленный цикл в капиталистической экономике представляет собой объективно обусловленную, исторически определенную и относительно периодическую форму движения расширенного воспроизводства совокупного общественного капитала, выражающуюся в закономерной последовательной смене кризиса, депрессии, оживления и подъема и возникающую вследствие накопления, обострения, временного насильственного разрешения и последующего воспроизводства внутренних противоречий капиталистического способа производства. Термин \textit{промышленный цикл} является наиболее содержательно точным для обозначения данного явления, поскольку подчеркивает его непосредственную связь с развитием капиталистической промышленности и позволяет разграничить его с иными формами циклического движения, в частности с аграрным, кредитным и финансовым циклами; при этом наиболее адекватную основу его изучения образуют марксистская методология и марксистская политическая экономия, рассматривающие цикл как форму движения общественного капитала, тогда как конкретные исторические, статистические, экономико-математические и иные инструменты анализа должны применяться во взаимосвязи и в соответствии с принципом восхождения от абстрактного к конкретному;
    \item Ключевым непосредственным индикатором промышленного цикла является индекс промышленного производства, наиболее последовательно отражающий изменения материального выпуска, загрузки производственных мощностей и общего состояния промышленного капитала. Цикл общенационального масштаба возникает лишь в развитой капиталистической экономике, характеризующейся господством крупного машинного производства и высокой долей основного постоянного капитала, а его историческое становление связано с кризисом общего перепроизводства 1825 г. в Великобритании; исходной сферой циклического движения выступает промышленное производство как материальная основа капиталистического хозяйства, фундаментальной причиной – имманентные противоречия капиталистического строя, а ключевым фактором периодичности – массовое обновление основного капитала, обостряющее данные противоречия. Продолжительность промышленного цикла определяется главным образом сроками физического и морального износа основного капитала, но не является жестко фиксированной и изменяется под воздействием технологических, кредитных, институциональных, внешнеэкономических и иных факторов; наиболее адекватно его внутреннюю структуру отражает четырехфазовая модель \textquotedbl{}кризис $\to$ депрессия $\to$ оживление $\to$ подъем\textquotedbl{}, позволяющая разграничить качественно различные состояния воспроизводства, причем, несмотря на существование в капиталистической экономике иных видов циклов, доминирующее положение занимает промышленный цикл, оказывающий определяющее воздействие на их динамику;
    \item Промышленный цикл представляет собой специфическую форму диалектического развития капиталистической экономики на ее определенном историческом этапе, в рамках которой ее главные противоречия остаются неизменными, а второстепенные противоречия частично и временно снимаются в ходе кризиса перепроизводства, и в дальнейшем воспроизводятся вновь на более высокой стадии развития системы. Вследствие этого кризис не устраняет основы капиталистического способа производства, а создает условия для нового расширения накопления, в ходе которого ранее разрешенные противоречия возникают вновь на более высокой ступени развития системы;
    \item Воспроизводство основного капитала выступает ключевым материальным фактором промышленного цикла, поскольку специфика его длительного функционирования, постепенного перенесения стоимости, физического и морального износа требует периодического массового возмещения машин, оборудования и иных элементов производственного аппарата, а соответствующие инвестиционные волны усиливают противоречия между расширением производительных возможностей и ограниченными условиями прибыльной реализации. В капиталистической экономике моменты авансирования и обновления основного капитала имеют тенденцию к синхронизации между различными предприятиями под воздействием межфирменной конкуренции, изменения рыночной конъюнктуры, динамики нормы прибыли и необходимости внедрения более производительной техники, тогда как возникающие отклонения отдельных капиталов от отраслевого ритма постепенно нивелируются теми же факторами и не нарушают общей периодичности массового воспроизводства основного капитала;
    \item При датировке периодов и фаз промышленного цикла необходимо строго разграничивать кризисы общего и частичного перепроизводства, а также циклические и промежуточные кризисы, поскольку их механическое объединение искажает действительную продолжительность и внутреннюю структуру цикла. Проведенный анализ экономики США показал, что средняя продолжительность промышленного цикла составила 120,1 месяца, то есть практически десять лет, при средней продолжительности кризиса 16,6 месяца, или 13,4 месяца с учетом промежуточных оживлений, депрессии – 3,3 месяца, оживления – 23 месяца, а подъема – 79,5 месяца, или 64,1 месяца с учетом промежуточных кризисов; таким образом, наиболее короткой фазой выступает депрессия, а наиболее продолжительной – подъем, причем сокращение депрессии с 6–8 до 1–2 месяцев может быть связано с развитием стимулирующих инструментов антикризисной политики. Согласно разработанному в исследовании сценарному прогнозу, следующий циклический кризис общего перепроизводства в экономике США может развернуться в 2028–2032 гг., а непосредственным поводом его открытого проявления способен стать растущий финансовый пузырь в сегменте искусственного интеллекта, который при этом следует рассматривать не как фундаментальную причину кризиса, а как возможный механизм его запуска и распространения;
    \item Проведенные автором расчеты на основе актуальных статистических данных подтверждают фундаментальные положения марксистской политической экономии о долгосрочной тенденции роста органического строения капитала и понижения нормы прибыли, а также показывают, что характер движения данных показателей закономерно изменяется на различных фазах промышленного цикла. В период подъема ускоренное накопление постоянного капитала и расширение инвестиций первоначально обеспечивают рост массы прибыли, однако по мере завершения цикла усиливают относительное вытеснение живого труда, перенакопление производственных мощностей и давление на норму прибыли, тогда как кризис посредством обесценения капитала и сокращения производства временно восстанавливает условия прибыльного накопления и тем самым создает материальные предпосылки для начала нового цикла.
\end{enumerate}

Таким образом, проведенное исследование позволило теоретически и эмпирически обосновать промышленный цикл как закономерную, исторически определенную и внутренне необходимую форму развития капиталистической экономики, коренящуюся в противоречиях воспроизводства общественного капитала и получающую материальную периодичность в процессе массового обновления основного капитала. Полученные результаты подтверждают, что государственное регулирование, кредитная экспансия, финансовые инновации и иные современные механизмы способны изменять продолжительность, глубину и внешние формы отдельных фаз цикла, однако не устраняют его фундаментальных причин, вследствие чего дальнейшее развитие капиталистического хозяйства неизбежно сохраняет циклический характер и периодически воспроизводит условия кризисов общего перепроизводства.
